\documentclass[10pt,a4paper]{amsart}
\usepackage[margin=19mm]{geometry}
\usepackage{amsmath,amssymb,mathtools}
\usepackage[scr=rsfs,cal=euler]{mathalfa}
\usepackage{xfrac}
\usepackage[unicode,hidelinks,bookmarks=false]{hyperref}
\newcommand{\ie}{i.e.\ }
\newcommand{\cf}{cf.\ }
\newcommand{\resp}{respectively\ }
\newcommand{\Subsection}{\subsection}
\providecommand{\nobreakdash}{\nobreak\hbox{-}}
\setlength{\emergencystretch}{2em}
\multlinegap0pt
\usepackage{amssymb}
\usepackage[all]{xy}
\newtheorem{lem}{Lemma}
\title{Generalized determinantal representation of hypersurfaces}
\author{A. El Mazouni, D. S. Nagaraj, Supravat Sarkar}
\date{Source: arXiv:2602.16685v1 (CC BY 4.0)}
\begin{document}
\maketitle

\noindent\emph{Source and licence.} Generalized determinantal representation of hypersurfaces. Version arXiv:2602.16685v1 (CC BY 4.0), distributed by arXiv under the Creative Commons Attribution 4.0 International licence, http://creativecommons.org/licenses/by/4.0/, as recorded in the arXiv licence metadata for this exact version and shown on the abstract page https://arxiv.org/abs/2602.16685v1. Full licence notice: \texttt{licenses/arxiv-2602.16685-CC-BY-4.0.txt}.\par\smallskip

\noindent\textit{Editorial scope.} Counted unit: the article's lem restriction (source0.tex lines 727--732). The article's complete original proof of it is delivered with it (source0.tex lines 733--753, one \texttt{proof} environment of 21 lines). No mathematics is written, repaired or re-derived: the statement and the proof are the article's own lines, in the article's own notation, and no retained line is substituted or re-flowed.  No further source-local context is needed: every cross-reference of every retained line resolves inside the two delivered spans.  Nothing was omitted from any retained span, so the package carries no omission record.  No retained line contains a citation, so the wrapper carries no bibliography and no bibliography entry.  No retained line contains a figure environment, an image-inclusion command or drawing code, so no image is delivered and the package declares no asset dependency.  Editorial formatting only, in addition: an AMS \texttt{amsart} preamble that restates the article's own declarations and macros used by the retained lines, namely the environments \texttt{lem} on separate counters, together with the standard packages those lines need.  The retained environments print in this document as Lemma 1 in order of appearance, whereas the article prints the same items with the numbering of its own full document.  The complete native source of the article as distributed in the arXiv e-print for this exact version is bundled unchanged as \texttt{sources/194912/source0.tex}.
\begin{lem}\label{restriction}
    Let $X$ a smooth projective surface with torsion free Picard group,
    $E$ a homogeneous vector bundle on $X$ which is $H$-abundant for 
    some ample line bundle $H$ on $X$. If $C$ is a smooth projective
    curve in $X$, then $E|_C$ is $H|_C$-abundant.
\end{lem}

\begin{proof}
    Let $d$ be the rank of $X$. For $m>>0$, we have a surjection 
    $$H^0(X,\textrm{det }E(mH))\to H^0(C,\textrm{det }E(mH)|_C).$$ Let
    $t\in H^0(C,\textrm{det }E(mH)|_C)$ be general. Choose $\tilde{t}\in 
    H^0(X,\textrm{det }E(mH)) $ mapping to $t$ under this surjection. As
    $t$ is general, we can assume $\tilde{t}$ is also general. As $E$ is
    $H$-abundant, there is $\phi\in \textrm{Aut }X$, an isomorphism
    \begin{equation}\label{isom}
        \phi^*\textrm{det }E(mH)\cong \textrm{det }E(mH),
    \end{equation}
and sections $\tilde{s_1},\tilde{s_2},...,\tilde{s_d}$ of $E(mH)$ 
    such that $\tilde{t}=\phi^*\tilde{s_1}\wedge 
    \phi^*\tilde{s_2}\wedge...\wedge \phi^*\tilde{s_d}$. As $E$ is
    homogeneous, \eqref{isom} shows $\phi^*(dmH)\cong dmH$, hence
    $dm(\phi^*H-H)=0$ in $\textrm{Pic }X$. As $\textrm{Pic }X$ is
    torison free, we get $\phi^*H\cong H$. Hence $\phi^*E(mH)\cong 
    E(mH)$. So we can regard $\phi^*\tilde{s_i}$ as a section of $
    E(mH)$, let $s_i$ be its restriction to $C$, hence a section of
    $E(mH)|_C$. Now $t$ is a nonzero scalar multiple of $s_1\wedge 
    s_2\wedge ...\wedge s_d$, completing the proof.
\end{proof}

\end{document}
